\begin{examplebox}
	\textbf{\textcolor{CExample}{例 1（基础）}}

	\textbf{\textcolor{CExample}{题目：}}
	已知圆内接四边形 $ABCD$ 中，$AB=3$，$BC=4$，$CD=5$，$DA=6$，求对角线 $AC$ 与 $BD$ 的乘积。

	\textbf{\textcolor{CExample}{解题步骤：}}
	\begin{description}[style=nextline, leftmargin=2.8em, labelsep=0.5em, itemsep=0.45em]
		\item[\textbf{第一步（识别条件）：}] 四边形 $ABCD$ 内接于圆，四边已知，符合定理条件。
			\par\textbf{理由：} 可直接使用托勒密定理。
		\item[\textbf{第二步（套用定理）：}] $AC\cdot BD = AB\cdot CD + AD\cdot BC$。
			\[
				\begin{aligned}
					AC\cdot BD &= AB\cdot CD + AD\cdot BC \\
					&= 3\times5 + 6\times4 \\
					&= 15+24 \\
					&= 39.
			\end{aligned}
			\]
			\par\textbf{理由：} 代入已知边长。
		\item[\textbf{第三步（得出结论）：}] 因此 $AC\cdot BD = 39$。
			\par\textbf{理由：} 计算完成。
	\end{description}

	\textbf{\textcolor{CExample}{关键结论：}}
	$\mathbf{AC\cdot BD = 39}$。

	\medskip
	\textbf{\textcolor{CExample}{例 2（稍变形）}}

	\textbf{\textcolor{CExample}{题目：}}
	圆内接四边形 $ABCD$ 中，$AB=2$，$BC=3$，$CD=4$，$DA=5$，且对角线 $AC=6$，求 $BD$ 的长。

	\textbf{\textcolor{CExample}{解题步骤：}}
	\begin{description}[style=nextline, leftmargin=2.8em, labelsep=0.5em, itemsep=0.45em]
		\item[\textbf{第一步（写出定理）：}] $AC\cdot BD = AB\cdot CD + AD\cdot BC$。
			\par\textbf{理由：} 四边形内接于圆，符合定理使用条件。
		\item[\textbf{第二步（代入数据）：}]
			\[
				\begin{aligned}
					6\cdot BD &= 2\times4 + 5\times3 \\
					&= 8+15 \\
					&= 23.
			\end{aligned}
			\]
			\par\textbf{理由：} 计算乘积和。
		\item[\textbf{第三步（解未知数）：}] $BD = \dfrac{23}{6}$。
			\par\textbf{理由：} 两边除以 $6$。
	\end{description}

	\textbf{\textcolor{CExample}{关键结论：}}
	$\mathbf{BD = \dfrac{23}{6}}$。

	\medskip
	\textbf{\textcolor{CExample}{例 3（综合应用）}}

	\textbf{\textcolor{CExample}{题目：}}
	在四边形 $ABCD$ 中，$AB=3$，$BC=4$，$CD=5$，$DA=6$，$AC\cdot BD = 39$。判断 $A,B,C,D$ 是否四点共圆。

	\textbf{\textcolor{CExample}{解题步骤：}}
	\begin{description}[style=nextline, leftmargin=2.8em, labelsep=0.5em, itemsep=0.45em]
		\item[\textbf{第一步（计算对边积和）：}]
			\[
				\begin{aligned}
					AB\cdot CD + AD\cdot BC &= 3\times5 + 6\times4 \\
					&= 15+24 \\
					&= 39.
			\end{aligned}
			\]
			\par\textbf{理由：} 直接计算。
		\item[\textbf{第二步（比较等式）：}] 已知 $AC\cdot BD=39$，恰好等于 $AB\cdot CD+AD\cdot BC$。
			\par\textbf{理由：} 两者相等。
		\item[\textbf{第三步（应用逆定理）：}] 由托勒密定理的逆定理，若 $AC\cdot BD = AB\cdot CD + AD\cdot BC$，则四边形 $ABCD$ 内接于圆。
			\par\textbf{理由：} 逆定理成立。
	\end{description}

	\textbf{\textcolor{CExample}{关键结论：}}
	$A,B,C,D$ 四点共圆。
\end{examplebox}
